\chapter{Information Metrics}

\section{Entropy}

\begin{definition}[Entropy]
    Let $X$ be a \textit{discrete} random variable with probability mass function $\PP_X(x)$, $x\in\cX$.
    \[H(X)=H(\PP_X)=\EE\left[\log\frac{1}{\PP_X(X)}\right]\]
\end{definition}
Here, we slightly abused the notation for a probability measure and its law.


\begin{definition}[Conditional entropy]
    \[H(X|Y)=\EE_{y\sim \PP_Y}[H(\PP_{X|Y=y})]=\EE\left[\log\frac{1}{\PP_{X|Y}(X|Y)}\right]\]
\end{definition}

\begin{example}
    Let $Y=X\oplus Z$, where $\oplus$ denotes binary addition (XOR) and $X\sim\text{Ber}(1/2)$ and $Z\sim\text{Ber}(\delta)$ independently of $X$.
    Then $P_{X|Y=0}=\text{Ber}(\delta)$ and $P_{X|Y=1}=\text{Ber}(\bar{\delta})$.
    Since $h(\delta)=h(\bar{\delta})$, we obtain that $H(X|Y)=h(\delta)$.
\end{example}

\begin{example}[Binary entropy function]
    \[h(p)=p\log\frac{1}{p}+\bar{p}\log\frac{1}{\bar{p}}\]
\end{example}

\begin{example}[Infinite Entropy]
    Let $\PP(X=k)\propto \frac{1}{k\ln^2k}$, $k=2,3,\dots$.
\end{example}

\begin{proposition}[Properties of entropy]
    \
    \begin{enumerate}[label=(\roman*)]
        \item $0\leq H(X)\leq \log|\cX|$
        \item $H(X|Y)\leq H(X)$
        \item $H(X,Y)=H(X)+H(Y|X)$
    \end{enumerate}
\end{proposition}

\subsection{Axiomatic Characterization}

\subsection{Submodularity}

\begin{example}[projection]
    
\end{example}

Let $S$ be a set. A set-function $f:2^S\to\RR$ is \textit{submodular} if for any $T_1,T_2\subset S$
\[f(T_1\cup T_2)+f(T_1\cap T_2)\leq f(T_1)+f(T_2).\]
Recall
\begin{theorem}[Submodularity]
    Let $X^n$ be a discrete random vector. Then $T\mapsto H(X_T)$ is submodular.
\end{theorem}
\begin{theorem}[Han's inequality]
    Let $X^n$ be a discrete $n$-dimensional random vector and denote $$\bar{H}_k(X^n)=\frac{1}{\binom{n}{k}}\sum_{|T|=k,T\subset[n]}H(X_T)$$
    the average entropy of a $k$-subset of coordinates. 
    Then, the sequence $\bar{H}_k$ is increasing and concave (in the sense of decreasing slope):
    \[\bar{H}_{k+1}-\bar{H}_k\leq\bar{H}_k-\bar{H}_{k-1}.\]
    As a consequence, $\frac{\bar{H}_k}{k}$ is decreasing in $k$.
\end{theorem}

Finally, we present a generaliation of Han's inequality.
\begin{theorem}[Shearer's Lemma]
    Let $X^n$ be a discrete $n$-dimensional random vector and let $S\subset [n]$ be a random variable independent of $X^n$. Then 
    \[H(X_S|S)\geq H(X^n)\cdot\min_{i\in [n]}\PP(i\in S).\]
\end{theorem}
\begin{proof}
    We prove an equivalent version: If $\cS=(S_1,\dots,S_M)$ is a list (possibly with repetitions) of subsets of $[n]$, then
    \[\sum_{j}H(X_{S_j})\geq H(X^n)\cdot\min_i\deg(i),\]
    where $\deg(i):=|\{j:i\in S_j\}|$. 
    
    If $\cS$ is a chain (all subsets can be rearranged so that $S_1\subset \dots\subset S_M$), then this statement is trivial,
    since the minimum on the RHS is either 0 (if $S_M\neq [n]$) or equals multiplicity of $S_M$ in $\cS$, in which case we have 
    \[\sum_j H(X_{S_j})\geq H(X_{S_M})\cdot |\{j:S_j=S_M\}|=H(X^n)\cdot\min_i\deg(i).\]
    
    When $\cS$ is not a chain, consider a pair of sets $S_1,S_2$ that are not related by inclusion and replace them in the collection with $S_1\cap S_2,S_1\cup S_2$.
    Submodularity implies that the LHS does not increase under this replacement, and the RHS is not changed.
    Notice that the total number of pairs that are not related by inclusion strictly decreases by this replacement.
    Therefore, by applying this operation we must eventually arrive to a chain, for which the statement has already been shown.
\end{proof}

Through the proof we can see that Han's inequality, 
Shearer's lemma holds for any submodular set-function that is also non-negative.

\begin{example}
    mutual information 
\end{example}

\section{Divergence}
\begin{definition}[Kullback-Leibler (KL) divergence]
    \[\KL(\PP\mmid \QQ)=\EE_{\QQ}\left[\frac{\d\PP}{\d\QQ}\log \frac{\d\PP}{\d\QQ}\right]\]
\end{definition}
Note that in general $\KL(\PP\mmid \QQ)\neq \KL(\QQ\mmid \PP)$. 
This is natural; for example, it reflects the inherent asymmetry of hypothesis testing.


for a pair of Markov kernels $\PP_{Y|X}:\cX\to\cY$
\begin{definition}[Conditional divergence]
    \[\KL(\PP_{Y|X}\mmid \QQ_{Y|X}\mid\PP_X)=\EE_{x\sim \PP_X}\KL(\PP_{Y|X=x}\mmid \QQ_{Y|X=x})\]
\end{definition}

\begin{example}[Binary divergence]
    Consider $\PP=\text{Ber}(p)$ and $\QQ=\text{Ber}(q)$ on $\cX=\{0,1\}$. Then 
    \[D(\PP\mmid\QQ)=d(p\mmid q):=p\log\frac{p}{q}+\bar{p}\log\frac{\bar{p}}{\bar{q}}.\]

    The following quadratic lower bound is often used:
    \[d(p\mmid q)\geq 2(p-q)^2\log e.\]
    In fact, this is the backbone of Pinsker's inequality.
\end{example}

\begin{proposition}
    \
    \begin{enumerate}[label=(\roman*)]
        \item $\KL(\PP\mmid\QQ)\geq 0$ with equality iff $\PP=\QQ$.
        \item \textnormal{(chain rule)} $\KL(\PP_{X,Y}\mmid\QQ_{X,Y})=\KL(\PP_{Y|X}\mmid \QQ_{Y|X}\mid\PP_X)+\KL(\PP_X\mmid \QQ_X)$
    \end{enumerate}
\end{proposition}

The chain rule has many implications

data-processing inequality (DPI)
the intuitive interpretation is that it is easier to distinguish two distributions using clean data as opposed to noisy data.

\subsection{Local Behavior and Fisher Information}

\subsection{Differential Entropy}

\section{Mutual Information}
We finally define the most important concept in the entire field of information theory, \textit{mutual information}.
\begin{definition}[Mutual information]
    The mutual information between a pair of random variables $X$ and $Y$ is
    \[I(X;Y)=\KL(\PP_{X,Y}\mmid \PP_X\PP_Y)\]
\end{definition}
Intuitively, mutual information measures the dependence between $X$ and $Y$ by comparing their joint distribution to the product of the marginals in the KL divergence.
Using the chain rule of KL-divergence, we see that mutual information can also be defined by comparing the conditional distribution to the unconditional:
\begin{proposition}
    Whenever $\cY$ is standard Borel,
    \[I(X;Y)=\KL(\PP_{Y|X}\mmid \PP_Y|\PP_X).\]
\end{proposition}

\begin{definition}[Conditional mutual information]
    \[I(X;Y|Z)=\EE_{z\sim P_Z}I(X;Y|Z=z).\]
\end{definition}


\begin{example}[Binary symmetric channel]
    Let $X\sim\text{Ber}(1/2)$ and $N\sim\text{Ber}(\delta)$ be independent.
    Let $Y=X\oplus N$; or equivalently, $Y$ is obtained by flipping $X$ with probability $\delta$.
    Then 
    \[I(X;Y)=\log 2 - h(\delta).\]
\end{example}



\begin{theorem}
    \
    \begin{enumerate}[label=(\roman*)]
        \item (positivity) $I(X;Y)\geq 0$, with equality $I(X;Y)=0$ iff $X\ind Y$.
        \item (symmetry) $I(X;Y)=I(Y;X)$.
        \item (deterministic maps) For any function $f$ we have $I(f(X);Y)\leq I(X;Y)$.
        If $f$ is one-to-one with a measurable inverse, then $I(f(X);Y)=I(X;Y)$.
        \item (chain rule) $I(X,Y;Z)=I(X;Z)+I(Y;Z|X).$
        \item (more data $\Longrightarrow$ more information) $I(X_1,X_2;Z)\geq I(X_1;Z)$.
        \item (DPI) If $X\to Y\to Z$, then $I(X;Z)\leq I(X;Y)$.
    \end{enumerate}
\end{theorem}

\begin{theorem}[Mutual information and entropy]
    \
    \begin{enumerate}[label=(\roman*)]
        \item If $X$ is discrete, then $I(X;Y)+H(X|Y)=H(X)$.
        \item If both $X$ and $Y$ are discrete, then $I(X;Y)+H(X,Y)=H(X)+H(Y)$
    \end{enumerate}
\end{theorem}



\begin{theorem}[Joint vs marginal mutual information]
    \
    \begin{enumerate}[label=(\roman*)]
        \item Memoryless channel: If $\PP_{Y^n|X^n}=\prod \PP_{Y_i|X_i}$, then 
        \[I(X^n;Y^n)\leq \sum_{i=1}^{n}I(X_i;Y_i),\]
        with equality iff $\PP_{Y^n}=\prod \PP_{Y_i}$.
        \item Memoryless source: If $X_1,\dots,X_n$ are independent, then \[I(X^n;Y)\geq \sum_{i=1}^{n}I(X_i;Y),\]
        with equality iff $\PP_{X^n|Y}=\prod \PP_{X_i|Y}$ $\PP_Y$-a.s. 
        Consequently,
    \end{enumerate}
\end{theorem}
\begin{proof}
    The easiest proof is by the variational principles.

    For (i), note that $I(X^n;Y^n)=\inf_{\QQ_{Y^n}}\KL(\PP_{Y^n|X}\mmid\QQ_{Y^n}\mid\PP_X)$ and take $\QQ_{Y^n}=\prod_{i=1}^{n}\PP_{Y_i}$.
    
    For (ii), note that $I(X^n;Y)=\sup_{\QQ_{X^n|Y}}\EE\log \frac{\QQ_{X^n|Y}}{\prod_{i=1}^{n}\PP_{X_i}}$ and take $\QQ_{X^n|Y}=\prod_{i=1}^{n}\PP_{X_i|Y}$. 
\end{proof}

\subsection{Data Processing and Sufficient Statistic}

\subsection{Fano's Inequality}

For the first time of this book, we connect informational quantities to operational concepts.
The key idea is to bound operational quantities involving arbitary complicated algorithms using informational quantities which can be evaluated.

\begin{proposition}
    Let $|\cX|=M<\infty$. Then for any $\hat{X}\ind X$,
    \[H(X)\leq F_M(\PP(X=\hat{X}))\]
    where $F_M(p)=h(p)+\bar{p}\log(M-1)$.
\end{proposition}
\begin{proof}
    Consider an auxillary distribution $\QQ_{X,\hat X}=\UU_X\PP_{\hat X}$, where $\UU_X$ is the uniform distribution on $\cX$.
    Then $\QQ(X=\hat X)=\frac{1}{M}$. Applying the DPI for KL divergence to the data processor $(X,\hat X)\mapsto 1\{X=\hat X\}$ yields 
    \[d(\PP(X=\hat{X})\mmid\frac{1}{M})\leq \KL(\PP_{X,\hat{X}}\mmid \QQ_{X,\hat{X}})=\log M-H(X).\]
\end{proof}
\begin{corollary}
    If $P_{\text{max}} =\max_{x\in\cX}\PP_X(x)$, then $$H(X)\leq F_M(P_{\text{max}}),$$ with equality iff $P_X=(P_{\text{max}},\frac{1-P_{\text{max}}}{M-1},\dots,\frac{1-P_{\text{max}}}{M-1})$.
\end{corollary}

\begin{theorem}[Fano's inequality]
    Let $X\to Y\to \hat{X}$ and $P_{\text{max}} =\max_{x\in\cX}\PP_X(x)$.
    Denote $P_e=\PP(X\neq\hat X)$, then 
    \[I(X;Y)\geq (1-P_e)\log \frac{1}{P_{\text{max}}}-h(P_e).\]
\end{theorem}
\begin{proof}
    
\end{proof}
Often, we further bound $h(P_e)\leq \log 2$ in statistical lower bound. 
This step is not very lossy.

\section{Variational Characterizations}

\subsection{Geometric Interpretation of Mutual Information}
\begin{theorem}[Golden formula]
    For any $\QQ_Y$ we have 
    \[\KL(\PP_{Y|X}\mmid \QQ_Y\mid \PP_X)=I(X;Y)+\KL(\PP_Y\mmid \QQ_Y).\]
    Thus, if $\KL(\PP_Y\mmid \QQ_Y)<\infty$, then 
    \[I(X;Y)=\KL(\PP_{Y|X}\mmid \QQ_Y\mid \PP_X)-\KL(\PP_Y\mmid \QQ_Y).\]
\end{theorem}
\begin{proof}
    According to the chain rule, we can decompose $\KL(\PP_{X,Y}\mmid \QQ_{X,Y})$, where $\QQ_{X,Y}=\PP_X\QQ_Y$, in two ways. 
    This leads to 
    \[\KL(\PP_X\mid\PP_X)+\KL(\PP_{Y|X}\mmid \QQ_Y\mid \PP_X)=\KL(\PP_Y\mmid\QQ_Y)+\KL(\PP_{X|Y}\mmid \QQ_{X|Y}\mid\PP_Y).\]
    Noting that $\QQ_{X|Y}=\PP_X$ completes the proof.
\end{proof}

\begin{corollary}[Mutual information as center of gravity]
    \[I(X;Y)=\min_{Q_Y}\KL(P_{Y|X}\mmid Q_Y\mid P_X).\]
\end{corollary}

\begin{corollary}[Mutual information as distance to product distribution]
    \[I(X;Y)=\min_{Q_X,Q_Y}\KL(P_{X,Y}\mmid Q_XQ_Y).\]
\end{corollary}

\begin{proposition}
    For any Markov kernel $Q_{X|Y}$ such that $Q_{X|Y=y}$ 
    If $I(X;Y)<\infty$, then 
    \[I(X;Y)=\sup_{Q_{X|Y}}\EE_{(X,Y)\sim \PP_{X,Y}}\log \frac{\d \QQ_{X|Y}}{\d \PP_X}\]
\end{proposition}
\begin{proof}
    By definition,
    \begin{align*}
        I(X;Y)&=\EE\log\frac{\d\PP_{X|Y}}{\d\PP_X}\frac{\d\QQ_{X|Y}}{\d\QQ_{X|Y}}\\
        &=\EE\log\frac{\d\QQ_{X|Y}}{\d\PP_X}+\EE\log\frac{\d\PP_{X|Y}}{\d\QQ_{X|Y}}\\
        &=\EE\log\frac{\d\QQ_{X|Y}}{\d\PP_X}+\KL(\PP_{X|Y}\mmid\QQ_{X|Y}\mid\PP_Y).
    \end{align*}
    $$$$
\end{proof}



\subsection{Characterizations of Divergence: Gelfand-Yaglom-Perez}
\begin{theorem}[Gelfand-Yaglom-Perez]
    Let $\PP,\QQ$ be two probability measures on $(\cX,\cF)$.
    Then 
    \[\KL(\PP\mmid\QQ)=\sup_{\{E_1,\dots,E_n\}}\sum_{i=1}^{n}\PP(E_i)\log\frac{\PP(E_i)}{\QQ(E_i)},\]
    where the supremum is over all finite $\cF$-measurable partitions.
\end{theorem}

\subsection{Characterizations of Divergence: Donsker-Varadhan}

For KL divergence, 

Defnie the tilting of $\QQ$ along $f$ by $\d{\QQ^f}=e^f\d{\QQ}$.
Now,
\begin{align*}
    \KL(\PP\mmid\QQ)&=\EE_P \log \frac{\d{\PP}}{\d{\QQ}}\\
    &=\EE_P \log \frac{\d{\PP}}{\d{\QQ^f}}\frac{\d{\QQ^f}}{\d{\QQ}}
\end{align*}

\begin{theorem}[Donsker-Varadhan]
    \[\]
\end{theorem}

\subsection{Gibbs Variational Principle}

\begin{theorem}[Gibbs variational principle]\label{thm:gibbs-variational-principle}
    \[\log \EE_\QQ e^f=\sup_{\PP} \left(\EE_\PP f-\KL(\PP\mmid\QQ)\right)\]
\end{theorem}



\section{Continuity}


\begin{proposition}
    Let $\cX$ be finite. Fix a distribution $\QQ>0$ on $\cX$.
    Then the map $\PP\mapsto \KL(\PP\mmid \QQ)$ is continuous.
\end{proposition}
\begin{remark}
    Divergence is not continuous in the pair, even for finite alphabets.
For example, as $n\to\infty$, $d(\frac{1}{n}\mmid 2^{-n})\not\to 0$.
\end{remark}
\begin{corollary}
    $P\mapsto H(P)$ is continous.
\end{corollary}

\begin{theorem}[Lower semicontinuity of divergence]
    Let $\cX$ be a metric space with Borel $\sigma$-algebra $\cF$.
    If $\PP_n$ and $\QQ_n$ converge weakly to $\PP$ and $\QQ$, respectively, then 
    \[\liminf_{n\to\infty} \KL(\PP_n\mmid\QQ_n)\geq \KL(\PP\mmid \QQ).\]
\end{theorem}
\begin{corollary}
    If $(X_n,Y_n)\stackrel{d}{\to}(X,Y)$, then 
    \[\liminf_{n\to\infty} I(X_n;Y_n)\geq I(X;Y).\]
\end{corollary}

\begin{proposition}[Continuity of mutual information]
    \
    \begin{enumerate}[label=(\roman*)]
        \item If $\cX$ and $\cY$ are both finite, then $P_{X,Y}\mapsto I(X;Y)$ is continuous.
        \item If $\cX$ is finite, then $P_X\mapsto I(X;Y)$ is continuous.
        \item Let $P_X$ range over the convex hull $\Pi =\co(P_1,\dots,P_n)$.
        If $I(P_j,P_{Y|X})<\infty$ for $j\in [n]$, then the map $P_X\mapsto I(X;Y)$ is continuous.
    \end{enumerate}
\end{proposition}

\section{Convexity}

\begin{theorem}[Entropy]
    The map $P_X\mapsto H(X)$ is concave.
    Furthermore, if $P_{Y|X}$ is any channel, then $\PP_X\mapsto H(X|Y)$ is concave.
    If $\cX$ is finite, then $P_X\mapsto H(X|Y)$ is continuous.
\end{theorem}

\begin{theorem}[KL divergence]
    The map $(\PP,\QQ)\mapsto \KL(\PP\mmid\QQ)$ is convex.
\end{theorem}

\begin{theorem}[Mutual information]
    \
    \begin{itemize}
        \item For fixed $P_{Y|X}$, $P_X\mapsto I(P_X,P_{Y|X})$ is concave.
        \item For fixed $P_X$, $P_{Y|X}\mapsto I(P_X,P_{Y|X})$ is convex.
    \end{itemize}
\end{theorem}

\section{Capacity: Saddle Point of Mutual Information}
Suppose we have a bivariate function $f$. Then we always have the minimax inequality:
\[\inf_y\sup_x f(x,y)\geq \sup_x\inf_y f(x,y).\]
If equality holds, we say that $f$ satisfies minimax theorem.


Recall that $(x^*,y^*)$ is a saddle point if 
\[f(x,y^*)\leq f(x^*,y^*)\leq f(x^*,y)\quad \forall x,y\]



\begin{lemma}
    
\end{lemma}

\begin{assumption}[Capacity-achieving input distribution]

\end{assumption}
\begin{example}[Non-existence of capacity-achieving input distribution]
    
\end{example}

It equals to the radius of the smallest divergence ``ball'' that encompasses all distributions $\{P_{Y|X=x}:x\in\cX\}$.
Moreover, the optimal center $P_Y^*$ is a convex combination of some $P_{Y|X=x}$ and is equidistant to those.

\begin{example}[Gaussian saddle point]
    
\end{example}

\begin{theorem}[\citet{kemperman1974shannon}]\label{thm:kemperman-capacity}
    For any $P_{Y|X}$ and a convex set of distributions $\cP$ such that
    \[C=\sup_{P_X\in\cP}I(P_X,P_{Y|X})<\infty,\]
    there exists a unique $P_Y^*$ with the property that 
    \[C=\sup_{P_X\in\cP}\KL(P_{Y|X}\mmid P_Y^*|P_X).\]
    Furthermore,
    \begin{align*}
        C&=\sup_{P_X\in\cP}\min_{Q_Y}\KL(P_{Y|X}\mmid Q_Y\mid P_X)\\
        &=\min_{Q_Y}\sup_{P_X\in\cP}\KL(P_{Y|X}\mmid Q_Y\mid P_X)
    \end{align*}
\end{theorem}

\section{Inequalities between f-divergences}



\section{f-Divergences}

\subsection{Definition}
\begin{definition}
    Let $f:(0,\infty)\to\RR$ be a convex function with $f(1)=0$.
    Let 
    If $\PP\ll\QQ$                        
\end{definition}

$f$-divergences satisfy the data-processing inequality

\begin{theorem}
    \[\dis_f(\PP_Y\mmid\QQ_Y)\leq \dis_f(\PP_X\mmid\QQ_X)\]
\end{theorem}

$f$-divergences satisfy two additional properties: conditioning increases divergence (CID) and convexity in pair.
\begin{theorem}
    
\end{theorem}

\begin{remark}
    More generally, we may call functional $g$-divergences.
\end{remark}

\subsection{Inequalities between f-Divergences: Joint Range}

\subsubsection{A Selection of Inequalities}

\begin{itemize}
    \item KL versus TV
    \item KL versus Hellinger
    \item KL versus $\chi^2$
    \item TV versus Hellinger
    \item TV versus $\chi^2$
\end{itemize}

\subsection{Variational Principle}

In general, 
\begin{theorem}
    $\dis_f(\PP\mmid\QQ)$
\end{theorem}



\section{Renyi Divergences}

The following family of divergence measures introduced in \cite{renyi1961measures} is key in many applications involving product measures.
Although these measures are not $f$-divergences, they are obtained as the monotone transformation of an appropriate $f$-divergence and thus satisfy DPI and other properties of $f$-divergences.

\begin{definition}
    For any $\lambda\in\RR\setminus\{0,1\}$, the R\'{e}nyi divergence of order $\lambda$ between probability distributions $\PP$ and $\QQ$ is defined as 
    \[\dis_\lambda(\PP\mmid\QQ):=\frac{1}{\lambda-1}\log\EE_\QQ \left[\left(\frac{\d\PP}{\d\QQ}\right)^\lambda\right],\]
    where $\EE_\QQ [(\frac{\d\PP}{\d\QQ})^\lambda]$ is formally understood as a $\sgn(\lambda-1)\dis_f(\PP\mmid\QQ)+1$ with $f(x)=\sgn(\lambda-1)x^\lambda-1$.
    The conditional R\'{e}nyi divergence is defined as 
\begin{align*}
    \dis_\lambda(\PP_{X|Y}\mmid\QQ_{X|Y}\mid \PP_Y)&:=\dis_\lambda (\PP_Y\times \PP_{X|Y}\mmid \PP_Y\times \QQ_{X|Y})\\
    &=\frac{1}{\lambda-1}\log\EE_{Y\sim\PP_Y}\int_\cX (\d\PP_{X|Y}(x))^\lambda(\d\QQ_{X|Y}(x))^{1-\lambda}.
\end{align*}
\end{definition}
\begin{example}
    Under mild regularity conditions $\lim_{\lambda\to 1}\dis_\lambda(\PP\mmid\QQ)=\KL(\PP\mmid\QQ)$.
    On the other hand, $\dis_2=\log(1+\chi^2)$ and $\dis_{\frac{1}{2}}=-2\log(1-\frac{H^2}{2})$ are monotone transformations of the $\chi^2$-divergence and the Hellinger distance, respectively.
\end{example}

\begin{proposition}
    \[\dis_\lambda\left(\prod_i\PP_{X_i}\mmid\prod_i\QQ_{X_i}\right)=\sum_i\dis_\lambda(\PP_{X_i}\mmid\QQ_{X_i}).\]
\end{proposition}
\begin{remark}
    The $\dis_\lambda$'s are the only divergences satisfying DPI and tensorization~\citep{mu2021blackwell}.
\end{remark}


\chapter{Hypothesis Testing and Large Deviations}

\section{Neyman-Pearson Lemma}

\subsection{Neyman-Pearson Formulation}
Suppose we have two given distribution $\PP$ and $\QQ$ on a space $\cX$.
We observe a sample $X$, which is generated from one of these distributions.
The two possibilities are summarized as a pair of hypotheses:
\begin{align*}
    &H_0:X\sim \PP, \\
    &H_1:X\sim \QQ.
\end{align*}
A (randomized) test is a Markov kernel $P_{Z|X}:\cX\to\{0,1\}$, so that $Z=0$ means the test choose $\PP$.

\begin{remark}
    This setting is called ``testing simple hypothesis against simple hypothesis.''
In comparison, a ``composite'' hypothesis postulates that $X\sim \PP$ for \textit{some} $\PP\in\cP$, where $\cP$ is a class of distributions.
\end{remark}

We can quantify the performance of a test by two performance metrics
\begin{equation*}
\begin{matrix}
 \alpha=\PP(Z=0) & 1-\text{type I error},\\
\beta=\QQ(Z=0)& \text{type II error}.
\end{matrix}
\end{equation*}
As we have two performance metrics, it is not easy to understand which test should be designated as the ``best''.
Consequently, there are several approaches:
\begin{itemize}
    \item Bayesian
    \item Minimax
    \item Neyman-Pearson: Minimize the type II error $\beta$ subject to the condition that the success probability under the null is at least $\alpha$.
\end{itemize}
To study the Neyman-Pearson formulation, we formalize the fundamental performance limit as follows:
\begin{definition}
    The Neyman-Pearson region consists of achievable points for all randomized tests
    \begin{equation*}
        \cR(\PP,\QQ)=\bigcup_{P_{Z|X}:\cX\to [0,1]}(\PP(Z=0),\QQ(Z=0))\subset [0,1]^2.
    \end{equation*}
    In particular, its lower boundary is defined as
    \begin{equation*}
        \beta_\alpha(\PP,\QQ)=\inf_{\PP(Z=0)\geq \alpha}\QQ(Z=0).
    \end{equation*}
\end{definition}
\begin{remark}
    The Neyman-Pearson region encodes much useful information about the relationship between $\PP$ and $\QQ$.
    In particular, every $f$-divergence can be computed from $\cR(\PP,\QQ)$.
\end{remark}

\section{Information Projection and Large Deviations}

\section{Large-Deviations Exponents}
\begin{theorem}
    Let $X_1,\dots,X_n\stackrel{i.i.d.}{\sim} \PP$.
    Then for any $\gamma\in\RR$,
    \begin{equation*}
        \lim_{n\to\infty}\frac{1}{n}\log \frac{1}{\PP(\frac{1}{n}\sum_{i=1}^{n}X_i>\gamma)}=\inf_{\QQ:\,\EE_\QQ X>\gamma}\KL(Q\mmid P)
    \end{equation*}
    Further, for every $n$ we have the firm upper bound 
    \begin{equation*}
        \PP\left(\frac{1}{n}\sum_{i=1}^{n}X_i>\gamma\right)\leq \exp\left\{-n\cdot \inf_{\QQ:\,\EE_\QQ X>\gamma}\KL(Q\mmid P)\right\}.
    \end{equation*}
\end{theorem}
\begin{proof}
    Let $E_n=\{ \frac{1}{n}\sum_{i=1}^{n}X_i>\gamma \}$.

    \textbf{Upper bound on $\PP(E_n)$.}

    \textbf{Lower bound on $\PP(E_n)$.} Let $\QQ$ be any distribution such that $\EE_\QQ X>\gamma$. Then by WLLN, $\QQ(E_n)=1-o(1)$.
    Now the DPI gives 
    \[d(\QQ(E_n)\mmid \PP(E_n))\leq \KL(\QQ_{X^n}\mmid \PP_{X^n})=n\KL(\PP\mmid \QQ).\]
    We also know that 
\end{proof}
\begin{corollary}
    The above result still holds for $\geq$ in place of $>$.
\end{corollary}

For a general set of distribution $\cE$
\begin{theorem}[Sanov's theorem, informal]
    Given $X_1,\dots,X_n\stackrel{i.i.d.}{\sim} \PP$ on $\cX$, denote the empirical distribution by $\hat{P}=\frac{1}{n}\sum_{j=1}^{n}\delta_{X_j}$.
    Let $\cE$ be a set of distributions. Then under regularity conditions on $\cE$ and $\PP$,
\begin{equation*}
    \PP(\hat{P}\in\cE)=\exp\left\{-n\min_{\QQ\in\cE}\KL(\QQ\mmid\PP)+o(n)\right\}
\end{equation*}
\end{theorem}
Regularity includes citation.

\section{Information Projection}


\section{Sequential Hypothesis Testing}

So far we have always been working with a fixed number of observations $n$.
However, different realizations of $X^n$ are informative to different 
In the sequential setting, pioneered by Wald \cite{}, the tester is allowed to request more observations.

Formally, a sequential hypothesis test consists of:
\begin{enumerate}
    \item a stopping time $\tau$ wit respect to the filtration $\{\cF_k,k\in\ZZ_+\}$, where $\cF_k:=\sigma(X_1,\dots,X_k)$ is generated by the first $k$ observations;
    \item a random variable (decision) $Z\in\{0,1\}$ measurable with respect to $\cF_\tau$.
\end{enumerate}
Each sequential test is associated with the following performance metrics:

\begin{example}
    The easiest way to see why sequential tests may be dramatically superior to fixed-sample-size tests is the following example.
    Consider $\PP=\frac{1}{2}\delta_0+\frac{1}{2}\delta_1$ and $\QQ=\frac{1}{2}\delta_0+\frac{1}{2}\delta_{-1}$.
    It is clear that no finite-sample-size test can achieve zero error under both hypotheses.
    However, an obvious sequential test (waiting for the first appearance of $\pm1$) achieves zero error probability, with a finite number of observations in expectation under both hypotheses.
\end{example}



\section{Composite Hypothesis Testing}

\begin{align*}
    &H_0:X\sim \PP, \quad \PP\in\cP\\
    &H_1:X\sim \QQ, \quad \QQ\in\cQ.
\end{align*}
where $\cP$ and $\cQ$ are two families of distributions.

Sometimes

uniform most powerful (UMP) tests

\begin{example}[Robust Hypothesis Testing]
    
\end{example}

\chapter{Lossless Data Compression}

\section{Universal Compression}

Since the probability law is ususally not given to us, one cannot apply the compression results that we have discussed so far.
In this section, we discuss how to produce compression schemes that do not require a priori knowledge of the distribution.
We will see that the optimal algorithms for this problem often coincide with the optimal algorithms for density estimation.

\subsection{Online Prediction and Density Estimation}

It turns out that the universal compression problem
automatically solves two other important problems: online prediction in machine learning and density estimation in statistics.


\chapter{Channel Coding}

\section{Shannon's Noisy Channel Coding Theorem}

\chapter{Rate Distortion Theory}

\section{Scalar and Vector Quantization}

\section{Shannon's Rate-Distortion Theorem}

\subsection{Information-Theoretic Formulation}
\paragraph{Hamming Game} 
Given 100 unbiased bits, we are asked to inspect them and scribble something down on a piece of paper that can store 50 bits at most.
Later we will be asked to guess the original 100 bits, with the goal of maximizing the number of correctly guessed bits. What is the best strategy?
Intuitively, we can store half of the bits and randomly guess the rest, which gives 25\% bit error rate (BER).
However, we will show that the optimal strategy amazingly achieves a BER of 11\%.

\paragraph{Mathematical Formulation} A lossy compressor

\subsection{Converse Bounds}
\begin{definition}[Rate-distortion function]
    \begin{equation}\label{eq:rate-distortion function}
        \phi_X(D)=\inf_{P_{Y|X}:\EE[d(X,Y)]\leq D }I(X,Y)
    \end{equation}
\end{definition}
\begin{proposition}[Properties of rate-distortion function]
    \
    \begin{enumerate}[label=(\roman*)]
        \item $\phi_X$ is convex, non-increasing
    \end{enumerate}
\end{proposition}
\begin{theorem}[General converse]
    Suppose $X\to W\to \hat{X}$, where $W\in [M]$ and $\EE[d(X,\hat{X})]\leq D$. Then 

\end{theorem}

\section{Eva}
\section{Metric Entropy}
\begin{definition}
    Let $(V,d)$ be a metric space and $\Theta\subset V$.
    \begin{itemize}
        \item We say $\{v_1,\dots,v_N\}\subset V$ is an $\epsilon$-covering of $\Theta$ if $\Theta\subset \cup_{i=1}^N B(v_i,\epsilon)$, where $B(v,\epsilon):=\{u\in V:d(u,v)\leq\epsilon\}$ is the closed ball of radius $\epsilon$ centered at $v$.
        \item We say $\{\theta_1,\dots,\theta_M\}\subset\Theta$ is an $\epsilon$-packing of $\Theta$ if $\min_{i\neq j}\|\theta_i-\theta_j\|>\epsilon$.
    \end{itemize}
\end{definition}

covering and packing number 
\begin{align*}
    &N(\Theta,d,\epsilon):=\min\{n:\exists \epsilon\text{ -covering of }\Theta\text{ of size }n\},\\
    &MN(\Theta,d,\epsilon):=\max\{m:\exists \epsilon\text{ -packing of }\Theta\text{ of size }m\}.
\end{align*}

\begin{theorem}[Kolmogorov-Tikhomirov]
    \[M(\Theta,d,2\epsilon)\leq N(\Theta,d,\epsilon)\leq M(\Theta,d,\epsilon).\]
\end{theorem}
\begin{remark}
    As we will see (for example, in \Cref{thm:volume-bound}), it is often easier to prove negative results (lower bounds on the minimal covering or upper bounds on the maximal packing) than positive results which require explicit construction.
    When use in conjunction with $N(\epsilon)\leq M(\epsilon)$, these turn into achievability statements, leading to many useful bounds on metric entropy.
    Revisiting the proof, we see that this logic actually corresponds to a greedy construction (increase the packing until no points can be added).
\end{remark}

\subsection{Finite-Dimensional Space and Volume Bound}

A commonly used method to bound metric entropy in finite dimensions is in terms of the volume ratio.
Consider the $d$-dimensional Euclidean space $V=\RR^d$ with metric given by an arbitary norm $d(x,y)=\|x-y\|$.
\begin{theorem}\label{thm:volume-bound}
    Let $\|\cdot\|$ be an arbitary norm on $\RR^d$ and $B=\{x\in\RR^d:\|x\|\leq 1\}$ the corresponding unit-norm ball.
    Then for any $\Theta\subset\RR^d$,
    \[\left(\frac{1}{\epsilon}\right)^d\frac{\vol(\Theta)}{\vol(B)}\stackrel{(\rm a)}{\leq} N(\Theta,\|\cdot\|,\epsilon)\leq M(\Theta,\|\cdot\|,\epsilon)\stackrel{(\rm b)}{\leq} \frac{\vol(\Theta+\frac{\epsilon}{2}B)}{\vol(\frac{\epsilon}{2}B)}\stackrel{(\rm c)}{\leq}\left(\frac{3}{\epsilon}\right)^d\frac{\vol(\Theta)}{\vol(B)},\]
    where $(\rm c)$ holds under the extra condition that $\Theta$ is convex and contains $\epsilon B$.
\end{theorem}


Next, we examine the discrete case of Hamming space.
\begin{theorem}
    For any integer $1\leq r\leq d-1$,
    \[\frac{2^d}{\sum_{i=0}^{r}{d\choose i}}\leq M(\FF_2^d,r)\leq \frac{2^d}{\sum_{i=0}^{\lfloor r/2\rfloor}{d\choose i}}\]
\end{theorem}

\begin{theorem}[Gilbert-Varshamov bound for Hamming spheres]\label{thm:Gilbert-Varshamov}
    Denote by $S^d_k:=\{x\in\FF_2^d:w_{\rm H}(x)=k\}$ the Hamming sphere of radius $0\leq k\leq d$. Then 
    \[M(S_k^d,r)\geq \frac{{d\choose{k}}}{\sum_{i=0}^{r}{d\choose{i}}}.\]
    In particular,
    \[\log M(S_k^d,\frac{k}{2})\geq \frac{k}{2}\log\frac{d}{2ek}.\]
\end{theorem}
\begin{proof}
    For the second bound, note that $\sum_{i=0}^{\frac{k}{2}}{d\choose{i}}\leq \exp(dh(\frac{k}{2s}))$.   
    Using the inequality $h(x)\leq x\log\frac{e}{x}$ and ${d\choose{k}}\geq (\frac{d}{k})^k$ 
\end{proof}

\subsection{Beyond the Volume Bound}

The volume bound

However, this result does not apply if the dimension $d$ is large and $\epsilon$ scales with $d$.
In fact, one expects that there is some critical threshold of $\epsilon$ depending on the dimension $d$, below which the exponential asymptotics is tight, and above which the covering number can grow polynomially in $d$.

\begin{example}
    Consider $M(B_1,\|\cdot\|_2,\epsilon)$, where $B_q$ denotes the unit $\ell_q$ ball with $1\leq q\leq \infty$.
    Let us start by applying the volume method in \Cref{thm:volume-bound}:
    \[\frac{\vol(B_1)}{\vol(\epsilon B_2)}\leq N(B_1,\|\cdot\|_2,\epsilon)\leq M(B_1,\|\cdot\|_2,\epsilon)\leq \frac{\vol(B_1+\frac{\epsilon}{2}B_2)}{\vol(\frac{\epsilon}{2}B_2)}.\]
    Applying the formula for the volume of a unit $\ell_q$-ball in $\RR^d$:
    \[\vol(B_q)=\frac{\left(2\Gamma(1+\frac{1}{q})\right)^d}{\Gamma(1+\frac{d}{q})},\]
    which yield, by Stirling's approximation, 
    \[\vol(B_q)^{\frac{1}{d}}\asymp d^{-\frac{1}{q}}.\]
    Then for some absolute constant $C$,
    \[M(B_1,\|\cdot\|_2,\epsilon)\leq\frac{\vol(B_1+\frac{\epsilon}{2}B_2)}{\vol(\frac{\epsilon}{2}B_2)}= \left(C(1+\frac{1}{\epsilon\sqrt{d}})\right)^d,\]
    as $B_2\subset\sqrt{d}B_1$ and $\vol(B_1+\frac{\epsilon}{2}B_2)\geq \max\{\vol(B_1),\vol(\frac{\epsilon}{2}B_2)\}$.
    On the other hand, for some absolute constant $c$,
    \[M(B_1,\|\cdot\|_2,\epsilon)\geq\left(\frac{1}{\epsilon}\right)^d \frac{\vol(B_1)}{\vol(\frac{\epsilon}{2}B_2)}=\left(\frac{c}{\epsilon\sqrt{d}}\right)^d.\]
    Overall, for $\epsilon\leq \frac{1}{\sqrt{d}}$, we have $M(B_1,\|\cdot\|_2,\epsilon)^{\frac{1}{d}}\asymp \frac{1}{\epsilon\sqrt{d}}$; however, the lower bound is exponential small in $d$ and the upper bound is exponential large in $d$ in the regime $\epsilon\gg\frac{1}{\sqrt{d}}$.
\end{example}

Geometrically, the intuition is that the volume of the $B_1$ ball is almost entirely concentrated near its vertices (the sparse axes).
When $\epsilon$ is large, the volume bound averages over the whole space, missing the geometry of the spikes.
The following results describes the complete behavior of this metric entropy, using methods that force the problem onto the discrete spikes.


\begin{theorem}
    For $\epsilon\in(0,1)$ and $d\in\NN$,
    \[\log M(B_1,\|\cdot\|_2,\epsilon)\asymp \begin{cases}
d\log\frac{e}{\epsilon^2d}, &\epsilon \leq \frac{1}{\sqrt{d} } ,  \\
\frac{1}{\epsilon^2}\log(e\epsilon^2 d),&\epsilon \geq \frac{1}{\sqrt{d} }  ,
\end{cases}\]
\end{theorem}


\begin{proof}
    We focus on $\frac{1}{\sqrt{d}}\leq \epsilon<1$. The other case follows from earlier volume calculations.

    For the upper bound, we construct an $\epsilon$-covering in $\ell_2$ by quantizing each coordinate.
    Fix some $\delta<1$. For each $\bm\theta\in B_1$, there exists $\bx\in (\delta\ZZ^d)\cap B_1$ such that $\|\bx-\bm\theta\|_\infty\leq \delta$.
    Then $\|\bx-\bm\theta\|_2^2\leq \|\bx-\bm\theta\|_1\|\bx-\bm\theta\|_\infty\leq 2\delta$. Picking $\delta=\frac{\epsilon^2}{2}$, $\bx$ forms an $\epsilon$-covering.
    Now, $\frac{\bx}{\delta}$ belongs to the set \[\mathcal{Z}=\{\bz\in\ZZ^d:\sum_{i=1}^{d}|z_i|\leq k\}\]
    with $k=\lfloor\frac{2}{\epsilon^2}\rfloor$. Note that each $\bz\in\mathcal{Z}$ has at most $k$ non-zeros.
    By enumerating, we have $|\mathcal{Z}|\leq 2^{k\wedge d}{d+k\choose{k}}$.

    For the lower bound, we construct a packing of $B_1$ based on a packing of the Hamming sphere.
    Fix some $1\leq k\leq d$ to be determined. Applying the Gilbert-Varshamov bound in \Cref{thm:Gilbert-Varshamov}, there exists a $\frac{k}{2}$-packing $\{\bx_1,\dots,\bx_M\}\subset \{\bx\in\{0,1\}^d:\sum_{i=1}^{d}x_i=k\}$ with $\log M\geq \frac{k}{2}\log\frac{d}{2ek}$.
    Set $\bm\theta_i=\frac{\bx_i}{k}$. Then $\bm\theta_i\in B_1$ and $\|\bm\theta_i-\bm\theta_j\|_2^2=\frac{1}{k^2}d_{\rm H}(\bx_i,\bx_j)\geq \frac{1}{2k}$ for all $i\neq j$.
    Choosing $k=\lfloor\frac{1}{\epsilon^2}\rfloor$, we conclude that $\{\bm\theta_1,\dots,\bm\theta_M\}$ is an $\frac{\epsilon}{2}$-packing of $B_1$ in $\|\cdot\|_2$.
\end{proof}
\begin{remark}
The above elementary proof can be adapted to give the following more general result: Let $1\leq p<q\leq\infty$.
For all $0<\epsilon<1$ and $d\in\NN$, 
\[\log M(B_p,\|\cdot\|_q,\epsilon)\asymp_{p,q} \begin{cases}
d\log\frac{e}{\epsilon^s d}, &\epsilon \leq d^{-\frac{1}{s}} ,  \\
\frac{1}{\epsilon^s}\log(e\epsilon^s d),&\epsilon \geq d^{-\frac{1}{s}}  ,
\end{cases}\quad \frac{1}{s}:=\frac{1}{p}-\frac{1}{q}.\]
\end{remark}
\begin{remark}
When $p>q$, we have 
\[\log M(B_p,\|\cdot\|_q,\epsilon)\asymp_{p,q} \begin{cases}
d\log\frac{e d^{\frac{1}{s}}}{\epsilon }, &\epsilon \leq d^{\frac{1}{s}} ,  \\
0,&\epsilon \geq d^{\frac{1}{s}}  ,
\end{cases}\quad \frac{1}{s}:=\frac{1}{q}-\frac{1}{p},\]
where the first line can be proved by the second line is because $B_p\subset d^{\frac{1}{s}}B_q$.
\end{remark}

\begin{remark}[Duality of Metric Entropy]
    \cite{artstein2004duality}
\end{remark}


\subsection{Infinite-Dimensional Space: Smooth Functions}

\begin{theorem}
    Assume that $L,A>0$ and $p\in [1,\infty]$ are constants.
    Consider the class $\cF(A,L)$ of all $L$-Lipschitz probability densities on the compact interval $[0,A]$.
    Then 
    \[\log N(\cF(A,L),\|\cdot\|_p,\epsilon)=\Theta(\frac{1}{\epsilon}).\]
    Furthermore, for $p=\infty$ we have the sharp asymptotics:
    \[\log_2 N(\cF(A,L),\|\cdot\|_\infty,\epsilon)=\frac{LA}{\epsilon}(1+o(1)),\quad\epsilon\to 0.\]
\end{theorem}
\begin{proof}
    By H\"{o}lder's inequality, $\|f\|_1\leq A^{1-\frac{1}{p}}\|f\|_p\leq A\|f\|_\infty$.
    So it suffices to prove the upper bound for $p=\infty$ and the lower bound for $p=1$.
    Specifically, we aim to show, by explicit construction,
    \begin{align*}
        &N(\cF,\|\cdot\|_\infty,\epsilon)\leq \frac{C}{\epsilon}2^{\frac{A}{\epsilon}},\\
        &M(\cF,\|\cdot\|_1,\epsilon)\geq 2^{\frac{c}{\epsilon}}.
    \end{align*}
    By replacing $f(x)$ by $\frac{1}{\sqrt{L}}f(\frac{x}{\sqrt{L}})$, we can consider the collection of $1$-Lipschitz densities on $[0,A]$.

    To construct a covering for $p=\infty$, set $x_k=k\epsilon$ for $k=1,\dots,n$. Here $n=\lceil\frac{A}{\epsilon}\rceil$.
    Note that as $f$ is a density, $f(0)\leq M$ for some constant $M$.
    Let $\cG$ be the collection of all lattice paths of $n$ steps starting from $(0,j\epsilon)$ for some $j\in\{0,\dots,J\}$, where $J=\lceil \frac{M}{\epsilon} \rceil$.
    In other word, each element of $\cG$ is a continuous piecewise linear function on each subinterval $I_k=[x_k,x_{k+1})$ with slope being either $+1$ or $-1$.
    Then $|\cG|\leq (J+1)2^n=O(\frac{1}{\epsilon}2^{\frac{A}{\epsilon}})$.
    It is also to see that $\cG$ is an $\epsilon$-covering by induction.

    To construct a packing for $p=1$, we perturb the uniform density.
    Define a function $T$ by 
    \[T(x)=\begin{cases}
        x1\{x\leq\epsilon\}+(2\epsilon-x)1\{x\geq\epsilon\}+\frac{1}{A}, &0\leq x\leq 2\epsilon, \\
        0, & \text{otherwise}.
    \end{cases}\]
    Let $n=\lceil\frac{A}{4\epsilon}\rceil$ and $a=2n\epsilon$. 
    For each $y\in\{0,1\}^n$, define a density $f_y$ on $[0,A]$ such that 
    \[f_y(x)=\begin{cases}
        \sum_{i=1}^{n}y_i T(x-2(i-1)\epsilon), &0\leq x\leq a,\\
        \text{linearly extend}, &a\leq x\leq A.
    \end{cases}\]
    Since $\int_0^a f_y=\frac{1}{2}+O(\epsilon)$, the slope of the linear extension is $O(\epsilon)$ for sufficiently small $\epsilon$, and $f_y\in \cF$.
    Further, we have $\|f_y-f_{y'}\|_1=d_{\rm H}(y,y')\|T\|_1=\epsilon^2d_{\rm H}(y,y')$.
    Invoking \Cref{thm:Gilbert-Varshamov}, we have an $\frac{n}{2}$-packing $\cY$ of $\{0,1\}^n$ with $|\cY|\geq 2^{cn}$.
    Thus $\{f_y:y\in\cY\}$ constitutes an $\frac{n\epsilon^2}{2}=\Theta(\epsilon)$-packing of $\cF$.

\end{proof}


\begin{theorem}
    Fix $L,A>0$, $p\in [1,\infty]$, and $d\in\NN$.
    Let $\beta>0$ and write $\beta=\ell+\alpha$, where $\ell\in\ZZ_+$ and $\alpha\in(0,1]$.
    Consider the class $\cF_\beta(A,L)$ of $\ell$-times continuously differentiable densities $f$ on $[0,A]^d$ whose $\ell$-th derivative is $(L,\alpha)$-H\"{o}lder continuous.
    Then 
    \[\log N(\cF_\beta(A,L),\|\cdot\|_p,\epsilon)=\Theta(\epsilon^{-\frac{d}{\beta}}).\]  
\end{theorem}

for the class of analytic functions, $\Theta((\log\frac{1}{\epsilon})^2)$.

\subsection{Metric Entropy and Small-Ball Probability}

The small-ball problem in probability theory concerns the behavior of the function 
\[\phi(\epsilon):=\log\frac{1}{\PP(\|X\|\leq \epsilon)}\]
as $\epsilon\to 0$, where $X$ is a random variable taking values on some real separable Banach space $(V,\|\cdot\|)$.


\begin{theorem}[\cite{kuelbs1993metric}]
    For all $\epsilon>0$,
    \[\phi(2\epsilon)-\log 2\leq \log N()\leq 2\phi(\frac{\epsilon}{2}).\]
\end{theorem}

\chapter{Statistical Applications}


\section{Mutual Information Method}

In this section, we describe a strategy for proving statistical lower bound called the \textit{mutual information method},
which entails comparing the amount of information data provides with the minimum amount of information needed to achieve a certain estimation accuracy.

We first fix some prior $\pi$ on $\Theta$ and we aim to lower bound the Bayes risk $R_\pi^*$ of estimating $\theta\sim \pi$ on the basis of $X$ with respect to some loss function $\ell$.
Let $\hat\theta$ be an estimator such that $\EE\ell(\theta,\hat\theta)\leq D$.
Then applying the DPI to the Markov chain $\theta\to X\to\hat\theta$, we have 
\[I(\theta;X)\geq I(\theta;\hat\theta)\geq \inf_{P_{\hat\theta|\theta}:\EE\ell(\theta,\hat\theta)\leq D}I(\theta;\hat\theta).\]
Note that:
\begin{itemize}
    \item The rightmost quantity can be interpreted as the minimum amount of information required to achieve a given estimation accuracy. This is precisely the rate-distortion function $\phi$ defined in~\eqref{eq:rate-distortion function}.
    \item The leftmost quantity can be interpreted as the amount of information provided by the data about the latent parameter. Sometimes it suffices to further upper bound it by the capacity of the channel $P_{X|\theta}$ by maximizing over all priors.
\end{itemize}
Therefore, we arrive at the following lower bound on the Bayes and hence the minimax risks:
\[R_\pi^*\geq \phi^{-1}(I(\theta;X))\geq \phi^{-1}(C).\]


``source-channel separation''

Three popular approaches for proving statistical lower bounds, namely, \textit{Le Cam's method}, \textit{Assouad's lemma}, and \textit{Fano's method}.
All three follow from the mutual information method, corresponding to different choice of prior $\pi$ for $\theta$, namely, the uniform distribution over a two-point set $\{\theta_0,\theta_1\}$, the hypercube $\{0,1\}^d$, and a packing.
While these methods are highly useful in determining the minimax rate for many problems, they are often loose with constant factors compared to the MIM.

\subsection{Application: Denoising a Vector in Gaussian Noise}

\subsection{Application: Denoising a Sparse Vector}

\subsection{Application: Community Detection}


\section{Lower Bounds by Reduction to Hypothesis Testing}
\subsection{Le Cam's Two-Point Method}
\begin{theorem}
    Suppose the loss function $\ell:\Theta\times\Theta\to\RR_{+}$ satisfies the following $\alpha$-triangle inequality for some $\alpha>0$: For all $\theta_0,\theta_1,\theta\in\Theta$,
    \[\ell(\theta_0,\theta_1)\leq \alpha(\ell(\theta_0,\theta)+\ell(\theta_1,\theta)).\]
    Then 
    \begin{equation*}
        \inf_{\hat\theta}\sup_{\theta\in\Theta}\EE_\theta\ell(\theta,\hat\theta)\geq \sup_{\theta_0,\theta_1\in\Theta}\frac{\ell(\theta_0,\theta_1)\ell(\theta_1,\theta_0)}{\alpha(\ell(\theta_0,\theta_1)+\ell(\theta_1,\theta_0))}(1-\TV(P_{\theta_0},P_{\theta_1})).
    \end{equation*}
\end{theorem}

\subsection{Assouad's Lemma}


\subsection{Fano's Method}

\begin{theorem}
    Let $d$ be a metric on $\Theta$. Fix an estimator $\hat\theta$.
    For any $T\subset \Theta$ and $\epsilon>0$,
    \begin{equation*}
        \PP\left(d(\theta,\hat\theta)\geq\frac{\epsilon}{2}\right)\geq 1-\frac{C(T)+\log 2}{\log M(T,d,\epsilon)},
    \end{equation*}
    where $C(T)=\sup I(\theta,X)$ is the capacity of the channel from $\theta$ to $X$ with input space $T$, with the supremeum taken over all distributions on $T$.
    Thus, for any $r>0$,
    \begin{equation*}
        \inf_{\hat\theta}\sup_{\theta\in\Theta}\EE_\theta d(\theta,\hat\theta)^r\geq \sup_{T\subset\Theta,\epsilon>0}\left(\frac{\epsilon}{2}\right)^r\left(1-\frac{C(T)+\log 2}{\log M(T,d,\epsilon)}\right).
    \end{equation*}
\end{theorem}
\begin{proof}
    Consider an $\epsilon$-packing $T'=\{\theta_1,\dots,\theta_M\}\subset T$. 
    Let $\theta$ be uniformly distributed on this packing.
    Given any estimator, construct a test by rounding $\hat\theta$ to $\tilde{\theta}=\arg\min_{\theta\in T'}d(\hat\theta,\theta)$.
    By the triangle inequality, $d(\theta,\tilde{\theta})\leq 2 d(\theta,\hat\theta)$.
    Thus, $\PP(\theta\neq\tilde{\theta})\leq \PP(d(\theta,\tilde{\theta})\geq \frac{\epsilon}{2})$.
    Then we can appeal to Fano's inequality 
    \[\PP(\theta\neq\tilde{\theta})\geq 1-\frac{I(\theta;X)+\log 2}{\log M(T,d,\epsilon)}.\]
    The proof is completed by $I(\theta;X)\leq C(T)$.
\end{proof}

When applying Fano's method, since it is often difficult to evaluate the capacity $C(T)$,
it is useful to 


\section{Entropic Bounds for Statistical Estimation} 

\subsection{Yang-Barron's Construction}

\[R^*_{\KL}(n):=\inf_{\hat P}\sup_{\theta\in\Theta}\EE_{P_\theta}\KL(P_\theta\mmid \hat{P})=\inf_{\hat P}\sup_{\theta\in\Theta}\int \KL(P_\theta\mmid\hat{P}(\cdot|X^n))P_\theta^{\otimes n}(\d x^n),\]
where the infimum is taken over all estimators $\hat{P}=\hat{P}(\cdot| X^n)$ which is a distribution on $\cX$.
\begin{theorem}\label{thm:Yang-Barron}
    Let $C_n$ denote the capacity of the channel $\theta\mapsto X^n\sim P_\theta^{\otimes n}$,
    where the supremum is over all priors of $\theta$ taking values in $\Theta$.
    Denote by 
     Then
    \begin{align*}
        R^*_{\KL}(n)&\leq \frac{C_{n+1}}{n+1}
    \end{align*}
    Conversely,
    \[\sum_{i=0}^n R^*_{\KL}(i)\geq C_{n+1}.\]
\end{theorem}

Any estimator, $\hat{P}=\hat{P}(\cdot|X^n)$, is a distribution on $\cX$ depending on $X^n$.
As such, $\hat{P}$ can be identified with a conditional distribution, say, $Q_{X_{n+1}|X^n}$, and we shall do so henceforth.
So let us introduce an unseen obsevation $X_{n+1}$ that is drawn from $P_\theta$ and independent of $X^n$ conditioned on $\theta$. 

The following lemma shows that the Bayes KL risk equals the conditional mutual information and the Bayes estimator is precisely $P_{X_{n+1}|X^n}$ (with respect to the joint distribution induced by the prior),
known as the \textit{predictive density estimator} in the statistics literature.
\begin{lemma}
    The Bayes risk for prior $\pi$ is given by
    \[R^*_{\KL,\pi}(n):=\inf_{\hat{P}} \int \KL(P_\theta\mmid\hat{P}(\cdot|X^n))P_\theta^{\otimes n}(\d x^n)\pi(\d\theta)=I(\theta;X_{n+1}|X^n),\]
    where $\theta\sim\pi$ and $(X_1,\dots,X_{n+1})\stackrel{\iid}{\sim}P_\theta$ conditioned on $\theta$.
    The Bayes estimator achieving this infimum is given by $\hat{P}(\cdot|x^n)=P_{X_{n+1}|X^n=x^n}$.
\end{lemma}
\begin{proof}
    We can calculate
    \begin{align*}
        \inf_{Q_{X_{n+1}|X^n}}\EE_{\theta,X^n}\KL(P_\theta\mmid Q_{X_{n+1}|X^n})
        &=\EE_{X^n}\inf_{Q_{X_{n+1}|X^n}}\EE_{\theta|X^n}\KL(P_{X_{n+1}|\theta}\mmid Q_{X_{n+1}|X^n})\\
        &=\EE_{X^n}\inf_{Q_{X_{n+1}|X^n}}\KL(P_{X_{n+1}|\theta,X^n}\mmid Q_{X_{n+1}|X^n}|P_{\theta|X^n})\\
        &=\EE_{X^n}\KL(P_{X_{n+1}|\theta,X^n}\mmid P_{X_{n+1}|X^n}|P_{\theta|X^n})\\
        &=I(\theta;X_{n+1}|X^n)
    \end{align*}
    where the first equality is by identifying $P_\theta$ with the conditional distribution 
    Note that we have the Markov chain $X^n\to\theta\to X_{n+1}$.
\end{proof}

Note that the operational meaning of $I(\theta;X_{n+1}|X^n)$ is the information provided by one extra observation about $\theta$ having already obtained $n$ observations.
In most cases, since $X^n$ will have already allowed $\theta$ to be consistently estimated as $n\to\infty$, the additional usefulness of $X_{n+1}$ is vanishing.
This is made precise by the following result.
\begin{lemma}
    $n\mapsto I(\theta;X_{n+1}|X^n)$ is a decreasing sequence.
\end{lemma}
\begin{proof}
    The easiest proof is to use the convexity of mutual information.
    Consider a sampling channel where the input is $\theta$ and the output is $X$ sampled from $P_\theta$.
    Let $I(\pi)$ denotes the mutual information when the input distribution is $\pi$, which is concave. Then
    $$I(\theta;X_{n+1}|X^n)=\EE_{X^n}I(\PP_{\theta|X^n})\leq \EE_{X^{n-1}}I(\PP_{\theta|X^{n-1}})=I(\theta;X_n|X^{n-1}).$$
\end{proof}

Together with the chain rule of mutual information, we get
$$I(\theta;X_{n+1}|X^n)\leq \frac{1}{n+1}I(\theta;X^{n+1}).$$

Now we are ready to prove Theorem~\ref{thm:Yang-Barron}.
\begin{proof}
    For the converse bound, first fix any prior $\pi$. Since the minimax risk dominates any Bayes risk, we have
    \begin{align*}
        \sum_{t=0}^{n}R^*_{\KL}(t)\geq \sum_{t=0}^{n}I(\theta;X_{t+1}|X^t)=I(\theta;X^{n+1}).
    \end{align*}    
    Optimizing over the prior completes the proof.

    For the upper bound, if we assume the minimax theorem holds so that $R^*_{\KL}=\sup_{\pi\in\cP(\Theta)}R^*_{\KL,\pi}$,
    then 

    Next, we present the original idea of \citet{yang1999information}, which does not assume the minimax theorem.
    The main idea is to consider a Bayes estimator but analyze it in the worst case.
\end{proof}

\begin{theorem}
    Let $\cQ=\{P_{B|A=a}:a\in\cA\}$ be a collection of distributions on some space $\cB$ and denote the capacity $C=\sup_{P_A\in\cP(\cA)}I(A;B)$.
    Then 
    \[C=\inf_{\epsilon>0}\{\epsilon^2+\log N_{KL}(\cQ,\epsilon)\},\]
    where 
    \[N_{\KL}(\cQ,\epsilon):=\min\left\{N:\exists Q_1,\cdots,Q_N \text{ s.t. }\forall a\in\cA,\exists i\in [N],\KL(P_{B|A=a}\mmid Q_i)\leq\epsilon^2\right\}\]
    is the KL covering number for the class $\cQ$.
\end{theorem}
\begin{proof}
    For the upper bound,

    For the lower bound, suppose $C<\infty$. Theorem~\ref{thm:kemperman-capacity}.
\end{proof}



\subsection{Pairwise Comparison \`{a} la Le Cam-Birg\'{e}}

\begin{theorem}[Le Cam--Birg\'{e}]
    Denote by $N_{\rm H}(\cP,\epsilon)$ the $\epsilon$-covering number of the set $\cP$ under the Hellinger distance.
    Let $\epsilon_n$ be such that 
    \[n\epsilon_n^2\geq \log N_{\rm H}(\cP,\epsilon) 1.\]
    Then there exists an estimator $\hat{P}=\hat{P}(X_1,\dots,X_n)$ such that for any $t\geq 1$,
    \[\sup_{P\in\cP}P(H(P,\hat{P})>4t\epsilon_n)\lesssim e^{-t^2}\]
    and, consequently, 
    \[\sup_{P\in\cP}\EE_P[H^2(P,\hat P)]\lesssim \epsilon_n^2.\]
\end{theorem}
\begin{proof}
    Abbreviate $\epsilon=\epsilon_n$ and $N=N_{\rm H}(\cP,\epsilon_n)$.
    Let $P_1,\dots,P_N$ be a maximal $\epsilon$-packing of $\cP$ under the Hellinger distance, which also serves as an $\epsilon$-covering.
\end{proof}

\subsection{Yatracos' Class and Minimum-Distance Estimator}  

\begin{theorem}[\citet{yatracos1985rates}]
    Let $X_1,\dots,X_n\stackrel{\iid}{\sim} P\in\cP$,
    where $\cP$ is a collection of distributions on a common measurable space $(\cX,\cE)$.
    For any $\epsilon>0$, there exists a proper estimator $\hat{P}=\hat{P}(X_1,\dots,X_n)\in\cP$, such that 
    \[\sup_{P\in\cP}\EE_{P}\TV(\hat{P},P)^2\leq C\left(\epsilon^2+\frac{\log N(\cP,\TV,\epsilon)}{n}\right).\]
\end{theorem}

For a loss function that is a distance, a natural idea for obtaining a proper estimator is the minimum-distance estimator
\[P_{\texttt{min-dist}}=\arg\min_{Q\in\cP}\TV(\hat{P}_n,Q),\]
where $\hat{P}_n=\frac{1}{n}\sum_{i=1}^{n}\delta_{X_i}$ is the empirical distribution.
However, this strategy does not make sense if the elements of $\cP$ have densities.
This is because TV, which compares two distributions over all measurable sets, is too strong.
The key idea is to replace TV by a proxy, which only inspects a ``low complexity'' family of sets.

To this end, let $\cA\subset\cE$ be a finite collection of measurable sets to be specified.
Define a pseudo-distance 
\[\texttt{dist}(P,Q)=\sup_{A\in\cA}|P(A)-Q(A)|.\]
It is easy to see that \texttt{dist} satisfis the triangle inequality. As a result, the estimator
\[\tilde{P}=\arg\min_{Q\in\cP}\texttt{dist}(Q,\hat{P}_n),\]
satisfies
\[\texttt{dist}(\tilde{P},P)\leq \texttt{dist}(\tilde{P},\hat{P}_n)+\texttt{dist}(\hat{P}_n,P)\leq 2\,\texttt{dist}(\hat{P}_n,P).\]

In addition, applying the binomial tail bound and the union bound, we have
\[\EE \texttt{dist}(P,\hat{P}_n)^2\leq\frac{C_0\log|\cA|}{n}\]
for some absolute constant $C_0$.

Now we need to choose $\cA$ such that $\texttt{dist}$ approximates the TV on $\cP$.
Consider an $\epsilon$-covering $\{Q_1,\cdots,Q_N\}$ of $\cP$ in TV. Define the set 
\[A_{ij}:=\left\{x:\frac{\d Q_i}{\d (Q_i+Q_j)}(x)\geq \frac{\d Q_j}{\d (Q_i+Q_j)}(x)\right\}\]
and the \textit{Yatracos class} 
\[\cA:=\{A_{ij}:i\neq j\in [N]\}.\]
This ensures that $\texttt{dist}(Q_i,Q_j)=\TV(Q_i,Q_j)$. Consequently,
\begin{lemma}
    \[\texttt{dist}(P,Q)\leq \TV(P,Q)\leq \texttt{dist}(P,Q)+4\epsilon\]
\end{lemma}
\begin{proof}
    For any $P,Q\in\cP$, there exist $i,j\in [N]$, such that $\TV(P,Q_i)\leq\epsilon$ and $\TV(Q,Q_j)\leq \epsilon$. Then we have
    \begin{align*}
        \TV(P,Q)&\leq \TV(P,Q_i)+\TV(Q_i,Q_j)+\TV(Q_j,Q)\\
                &\leq 2\epsilon + \texttt{dist}(Q_i,Q_j)\\
                &\leq 2\epsilon + \texttt{dist}(Q_i,P) +\texttt{dist}(P,Q)  + \texttt{dist}(Q,Q_j)\\
                &\leq 4\epsilon +\texttt{dist}(P,Q).
    \end{align*}
\end{proof}
Finally, 
\begin{align*}
    \TV(\tilde{P},P)&\leq \texttt{dist}(\tilde{P},P)+4\epsilon\\
                    &\leq 2\, \texttt{dist}(P,\hat{P}_n)+4\epsilon.
\end{align*}
Squaring both sides, we have
\[\EE\TV(\tilde{P},P)^2\leq 32\epsilon^2 + \frac{16C_0\log|N|}{n}.\]

\begin{remark}
    Yatracos' scheme works even if the model is mis-specified.
\end{remark}

\subsection{Application: Density Estimation over H\"{o}lder Classes}

\begin{theorem}\label{thm:density-estimation-Holder-l2}
    Given $X_1,\dots,X_n\stackrel{\iid}{\sim}f\in \cF$, the minimax $L^2$ risk over $\cF$ satisfies
    \[R^*_{L^2}(n;\cF):=\inf_{\hat f}\sup_{f\in\cF}\EE\|f-\hat{f}\|_2^2\asymp n^{-2\beta/2\beta+d}\]
\end{theorem}

We will prove this result by applying the entropic upper bound and the minimax lower bound based on Fano's inequality.
To convert KL risk into $L^2$,
\begin{lemma}
    Suppose both $f=\frac{\d\PP}{\d\mu}$ and $g=\frac{\d\QQ}{\d\mu}$ are upper- and lower-bounded by absolute constants $c$ and $C$ respectively.
    Then
    \[\frac{1}{4C}\int (f-g)^2\d\mu\leq H^2(\PP,\QQ)\leq \KL(\PP\mmid\QQ)\leq \chi^2(\PP\mmid\QQ)\leq \frac{1}{c}\int (f-g)^2\d\mu.\]
\end{lemma}
\begin{proof}
    For the upper bound, $\KL(\PP\mmid\QQ)\leq \chi^2(\PP\mmid\QQ)=\int\frac{(f-g)^2}{g}\d\mu\leq \frac{1}{c}\int (f-g)^2 \d\mu$.
    
    For the lower bound, $\KL(\PP\mmid\QQ)\geq H^2(\PP,\QQ)=\int \frac{(f-g)^2}{(\sqrt{f}+\sqrt{g})^2}\geq \frac{1}{4C}\int (f-g)^2\d\mu$.
\end{proof}

\begin{lemma}
    Let $\cF'$ denote the collection of $f\in\cF$ which is bounded from below by $\frac{1}{2}$. Then 
    \[R^*_{L^2}(n;\cF')\leq R^*_{L^2}(n;\cF)\leq 4R^*_{L^2}(n;\cF').\]
\end{lemma}
\begin{proof}
    The first inequlaity follows because $\cF'\subset\cF$.

    For the second inequality, fix some $f\in\cF$ and we observe $X_1,\dots,X_n\stackrel{\iid}{\sim} f$.
    Let $U_1,\dots,U_n$ be independently sampled from the uniform distribution on $[0,1]^d$.
    Define 
    \[Z_i=\begin{cases}
        U_i & \text{w.p.}\frac{1}{2},\\
        X_i & \text{w.p.}\frac{1}{2}.
    \end{cases}\]
    Then $Z_1,\dots,Z_n \stackrel{\iid}{\sim} g:=\frac{1}{2}(1+f)\in\cF'$. 
    Let $\hat{g}$ be an estimator that achieves the minimax risk $R^*_{L^2}(n;\cF')$ on $\cF'$/
    Consider the estimator $\hat{f}=2\hat{g}-1$, then $\|f-\hat{f}\|_2^2=4\|g-\hat{g}\|_2^2$.
    Taking the supremum over $f\in\cF$ proves the desired result.
\end{proof}

Now we can prove \Cref{thm:density-estimation-Holder-l2}.
\begin{proof}
    Using Yang-Barron's construction, we have
    \begin{align*}
        R^*_{L^2}(n;\cF')&\asymp R^*_{\KL}(n;\cF')\\
        &\lesssim \inf_{\epsilon>0}\{\epsilon^2+\frac{1}{n}\log N_{\KL}(\cF',\epsilon)\}\\
        &\asymp \inf_{\epsilon>0}\{\epsilon^2+\frac{1}{n}\log N(\cF',\|\cdot\|_2,\epsilon)\}\\
        &\asymp \inf_{\epsilon>0}\{\epsilon^2+\frac{1}{n\epsilon^{\frac{d}{\beta}}}\}\asymp n^{-\frac{2\beta}{2\beta+d}}.
    \end{align*}

    For the lower bound, we can apply Fano's inequality.
\end{proof}